\documentclass[10pt,a4paper]{amsart}
\usepackage{amsmath,amssymb,amsthm,graphicx}
% Retained minimal preamble for the Stacks Project chapter varieties.tex.
% Verbatim excerpt of the chapter preamble: Xy-pic setup (native preamble lines 14-18), the native theorem
% declarations with the mathematical macros (native preamble lines 154-217), and the project own no-op
% definitions of the reference/slogan/history display environments (native wrapper.tex lines 1-3) that the
% counted statement uses for its slogan. Copied mechanically; nothing added, removed or redefined.
% For commutative diagrams we use Xy-pic
\usepackage[all]{xy}

% We use 2cell for 2-commutative diagrams.
\xyoption{2cell}
% Theorem environments.
%
\theoremstyle{plain}
\newtheorem{theorem}[subsection]{Theorem}
\newtheorem{proposition}[subsection]{Proposition}
\newtheorem{lemma}[subsection]{Lemma}

\theoremstyle{definition}
\newtheorem{definition}[subsection]{Definition}
\newtheorem{example}[subsection]{Example}
\newtheorem{exercise}[subsection]{Exercise}
\newtheorem{situation}[subsection]{Situation}

\theoremstyle{remark}
\newtheorem{remark}[subsection]{Remark}
\newtheorem{remarks}[subsection]{Remarks}

\numberwithin{equation}{subsection}

% Macros
%
\def\lim{\mathop{\mathrm{lim}}\nolimits}
\def\colim{\mathop{\mathrm{colim}}\nolimits}
\def\Spec{\mathop{\mathrm{Spec}}}
\def\Hom{\mathop{\mathrm{Hom}}\nolimits}
\def\Ext{\mathop{\mathrm{Ext}}\nolimits}
\def\SheafHom{\mathop{\mathcal{H}\!\mathit{om}}\nolimits}
\def\SheafExt{\mathop{\mathcal{E}\!\mathit{xt}}\nolimits}
\def\Sch{\mathit{Sch}}
\def\Mor{\mathop{\mathrm{Mor}}\nolimits}
\def\Ob{\mathop{\mathrm{Ob}}\nolimits}
\def\Sh{\mathop{\mathit{Sh}}\nolimits}
\def\NL{\mathop{N\!L}\nolimits}
\def\CH{\mathop{\mathrm{CH}}\nolimits}
\def\proetale{{pro\text{-}\acute{e}tale}}
\def\etale{{\acute{e}tale}}
\def\QCoh{\mathit{QCoh}}
\def\Ker{\mathop{\mathrm{Ker}}}
\def\Im{\mathop{\mathrm{Im}}}
\def\Coker{\mathop{\mathrm{Coker}}}
\def\Coim{\mathop{\mathrm{Coim}}}

% Boxtimes
%
\DeclareMathSymbol{\boxtimes}{\mathbin}{AMSa}{"02}

%
% Macros for moduli stacks/spaces
%
\def\QCohstack{\mathcal{QC}\!\mathit{oh}}
\def\Cohstack{\mathcal{C}\!\mathit{oh}}
\def\Spacesstack{\mathcal{S}\!\mathit{paces}}
\def\Quotfunctor{\mathrm{Quot}}
\def\Hilbfunctor{\mathrm{Hilb}}
\def\Curvesstack{\mathcal{C}\!\mathit{urves}}
\def\Polarizedstack{\mathcal{P}\!\mathit{olarized}}
\def\Complexesstack{\mathcal{C}\!\mathit{omplexes}}
% \Pic is the operator that assigns to X its picard group, usage \Pic(X)
% \Picardstack_{X/B} denotes the Picard stack of X over B
% \Picardfunctor_{X/B} denotes the Picard functor of X over B
\def\Pic{\mathop{\mathrm{Pic}}\nolimits}
\def\Picardstack{\mathcal{P}\!\mathit{ic}}
\def\Picardfunctor{\mathrm{Pic}}
\def\Deformationcategory{\mathcal{D}\!\mathit{ef}}
\newenvironment{reference}{}{}
\newenvironment{slogan}{}{}
\newenvironment{history}{}{}
\usepackage{cleveref}
\crefname{originalreference}{original reference}{original references}
\Crefname{originalreference}{Original reference}{Original references}
\makeatletter
\newlabel{situation-glue}{{09MY}{source}{}{sourcepointer-1}{}}
\newlabel{situation-glue@cref}{{[originalreference][0][]09MY}{[1][0][]source}{}{}{}}
\makeatother
\crefname{originalreference}{original reference}{original references}
\Crefname{originalreference}{Original reference}{Original references}
\makeatletter
\newlabel{algebra-lemma-characterize-Dedekind}{{034X}{source}{}{sourcepointer-2}{}}
\newlabel{algebra-lemma-characterize-Dedekind@cref}{{[originalreference][0][]034X}{[1][0][]source}{}{}{}}
\makeatother
\crefname{originallemma}{original lemma}{original lemmas}
\Crefname{originallemma}{Original lemma}{Original lemmas}
\makeatletter
\newlabel{lemma-semi-local}{{09N1}{source}{}{sourcepointer-3}{}}
\newlabel{lemma-semi-local@cref}{{[originallemma][0][]09N1}{[1][0][]source}{}{}{}}
\makeatother
\crefname{originalreference}{original reference}{original references}
\Crefname{originalreference}{Original reference}{Original references}
\makeatletter
\newlabel{algebra-lemma-normal-domain-intersection-localizations-height-1}{{031T}{source}{}{sourcepointer-4}{}}
\newlabel{algebra-lemma-normal-domain-intersection-localizations-height-1@cref}{{[originalreference][0][]031T}{[1][0][]source}{}{}{}}
\makeatother
\crefname{originallemma}{original lemma}{original lemmas}
\Crefname{originallemma}{Original lemma}{Original lemmas}
\makeatletter
\newlabel{lemma-semi-local-dimension-one-conductor}{{09N2}{source}{}{sourcepointer-5}{}}
\newlabel{lemma-semi-local-dimension-one-conductor@cref}{{[originallemma][0][]09N2}{[1][0][]source}{}{}{}}
\makeatother
\crefname{originalreference}{original reference}{original references}
\Crefname{originalreference}{Original reference}{Original references}
\makeatletter
\newlabel{algebra-lemma-chinese-remainder}{{00DT}{source}{}{sourcepointer-6}{}}
\newlabel{algebra-lemma-chinese-remainder@cref}{{[originalreference][0][]00DT}{[1][0][]source}{}{}{}}
\makeatother
\title{Two-sided glueing of one-dimensional Noetherian semi-local domains whose tensor product maps onto the common field}
\author{The Stacks Project Authors}
\date{stacks-project:varieties.tex:local-snapshot-20261011; GNU FDL 1.2 or later}
\begin{document}
\maketitle
\noindent\textit{Editorial note.} All delivered mathematical text is the author's own, copied verbatim from the pinned Stacks Project chapter file varieties.tex; no statement or proof step was rewritten, reordered or removed.\par\medskip
\noindent\textit{Editorial note.} The counted span cites seven cross-chapter labels (including a situation), each recorded in assembly.external\_references with its canonical tag as a navigation pointer.\par\medskip
\noindent\textit{Editorial note.} The proof cites no bibliography item inside the counted span.\par\medskip
\noindent\textit{Allowed standard prerequisites (editorial).}\begin{itemize}
\item Characterization of Dedekind domains by the discrete valuation property of localizations. Stacks Project, Algebra, Lemma algebra-lemma-characterize-Dedekind (tag 034X). Gives that the localizations of the integrally closed B at maximal ideals are discrete valuation rings.
\item The glueing situation for one-dimensional semi-local rings. Stacks Project, Varieties, Situation situation-glue (tag 09MY). Supplies the ambient hypotheses relating A, B, K and R.
\end{itemize}\medskip
\ifdefined\hypertarget\hypertarget{sourcepointer-1}{}\fi
\noindent\textit{Source reference.} Original reference 09MY: sources/stacks-chow/tags.txt, source lines 12196--12196. Cross-chapter standard prerequisite cited by the author inside the counted span; it is used as a pointer and its own proof is not part of this item.\par\medskip
\ifdefined\hypertarget\hypertarget{sourcepointer-2}{}\fi
\noindent\textit{Source reference.} Original reference 034X: sources/stacks-chow/tags.txt, source lines 4108--4108. Cross-chapter standard prerequisite cited by the author inside the counted span; it is used as a pointer and its own proof is not part of this item.\par\medskip
\ifdefined\hypertarget\hypertarget{sourcepointer-3}{}\fi
\noindent\textit{Source reference.} Original lemma 09N1: sources/stacks-chow/tags.txt, source lines 12199--12199. Cross-chapter standard prerequisite cited by the author inside the counted span; it is used as a pointer and its own proof is not part of this item.\par\medskip
\ifdefined\hypertarget\hypertarget{sourcepointer-4}{}\fi
\noindent\textit{Source reference.} Original reference 031T: sources/stacks-chow/tags.txt, source lines 3996--3996. Cross-chapter standard prerequisite cited by the author inside the counted span; it is used as a pointer and its own proof is not part of this item.\par\medskip
\ifdefined\hypertarget\hypertarget{sourcepointer-5}{}\fi
\noindent\textit{Source reference.} Original lemma 09N2: sources/stacks-chow/tags.txt, source lines 12200--12200. Cross-chapter standard prerequisite cited by the author inside the counted span; it is used as a pointer and its own proof is not part of this item.\par\medskip
\ifdefined\hypertarget\hypertarget{sourcepointer-6}{}\fi
\noindent\textit{Source reference.} Original reference 00DT: sources/stacks-chow/tags.txt, source lines 509--509. Cross-chapter standard prerequisite cited by the author inside the counted span; it is used as a pointer and its own proof is not part of this item.\par\medskip
\begin{lemma}
\label{lemma-semi-local-both-side}
In Situation \ref{situation-glue} assume
\begin{enumerate}
\item $A$ is a Noetherian semi-local domain of dimension $1$,
\item $B$ is a Noetherian semi-local domain of dimension $1$,
\item $A \otimes B \to K$ is surjective.
\end{enumerate}
Then $\Spec(A) \to \Spec(R)$ and $\Spec(B) \to \Spec(R)$
are open immersions covering $\Spec(R)$ and $K = A \otimes_R B$.
\end{lemma}

\begin{proof}
Special case: $B$ is integrally closed in $K$. This means that
$B$ is a Dedekind domain
(Algebra, Lemma \ref{algebra-lemma-characterize-Dedekind})
whence all of its localizations at maximal ideals are discrete valuation rings.
Let $\mathfrak m_1, \ldots, \mathfrak m_r$ be the maximal ideals of
$B$. We set
$$
R_1 = A \times_K B_{\mathfrak m_1}
$$
Observing that $A \otimes_{R_1} B_{\mathfrak m_1} \to K$ is surjective
we conclude from Lemma \ref{lemma-semi-local} that $A$ and
$B_{\mathfrak m_1}$ define open subschemes covering $\Spec(R_1)$ and that
$K = A \otimes_{R_1} B_{\mathfrak m_1}$. In particular $R_1$ is a
semi-local Noetherian ring of dimension $1$. By induction we define
$$
R_{i + 1} = R_i \times_K B_{\mathfrak m_{i + 1}}
$$
for $i = 1, \ldots, r - 1$. Observe that $R = R_r$ because
$B = B_{\mathfrak m_1} \cap \ldots \cap B_{\mathfrak m_r}$ (see
Algebra, Lemma
\ref{algebra-lemma-normal-domain-intersection-localizations-height-1}).
It follows from the inductive procedure that $R \to A$ defines an
open immersion $\Spec(A) \to \Spec(R)$. On the other hand, the
maximal ideals $\mathfrak n_i$ of $R$ not in this open correspond
to the maximal ideals $\mathfrak m_i$ of $B$ and in fact
the ring map $R \to B$ defines an isomorphisms
$R_{\mathfrak n_i} \to B_{\mathfrak m_i}$ (details omitted; hint:
in each step we added exactly one maximal ideal to $\Spec(R_i)$).
It follows that $\Spec(B) \to \Spec(R)$ is an open immersion
as desired.

\medskip\noindent
General case. Let $B' \subset K$ be the integral closure of $B$.
See Lemma \ref{lemma-semi-local-dimension-one-conductor}.
Then the special case applies to $R' = A \times_K B'$.
Pick $x \in R'$ which is not contained in the maximal
ideals of $A$ and is contained in the maximal ideals of $B'$
(see Algebra, Lemma \ref{algebra-lemma-chinese-remainder}).
By Lemma \ref{lemma-semi-local-dimension-one-conductor}
there exists an integer $n$ such that $x^n \in R = A \times_K B$.
Replace $x$ by $x^n$ so $x \in R$. For every $y \in R'$ there exists
an integer $n$ such that $x^n y \in R$. On the other hand,
it is clear that $R'_x = A$. Thus $R_x = A$.
Exchanging the roles of $A$ and $B$ we also find an $y \in R$
such that $B = R_y$. Note that inverting both $x$ and $y$
leaves no primes except $(0)$. Thus $K = R_{xy} = R_x \otimes_R R_y$.
This finishes the proof.
\end{proof}

\end{document}
